\documentclass[12pt,a4paper]{report}

\usepackage[utf8]{inputenc}
\usepackage[T1]{fontenc}
\usepackage[french]{babel}
\usepackage[a4paper,margin=2.5cm]{geometry}
\usepackage{graphicx}
\usepackage{array}
\usepackage{tabularx}
\usepackage{longtable}
\usepackage{booktabs}
\usepackage{enumitem}
\usepackage{setspace}
\usepackage{xcolor}
\usepackage{hyperref}
\usepackage{tcolorbox}
\usepackage{fancyhdr}
\usepackage{titlesec}
\usepackage{float}

\hypersetup{
  colorlinks=true,
  linkcolor=teal!50!black,
  urlcolor=teal!50!black,
  citecolor=teal!50!black,
  pdftitle={Rapport PFE - E-Joutia},
  pdfauthor={El Bahlouli Adil, El Boualiti Oussama}
}

\definecolor{ejoutia}{HTML}{18B7AA}
\definecolor{ink}{HTML}{0F172A}
\definecolor{softgray}{HTML}{607082}

\pagestyle{fancy}
\fancyhf{}
\fancyhead[L]{\textit{E-Joutia}}
\fancyhead[R]{\textit{Rapport PFE}}
\fancyfoot[C]{\thepage}
\setlength{\headheight}{15pt}

\fancypagestyle{emptyclean}{
  \fancyhf{}
  \renewcommand{\headrulewidth}{0pt}
  \renewcommand{\footrulewidth}{0pt}
}

\titleformat{\chapter}[display]
  {\bfseries\Huge\color{ink}}
  {\filright\Large\chaptername\ \thechapter}
  {0.8ex}
  {\filright}

\setlist[itemize]{label=--,leftmargin=1.2cm}
\setlist[enumerate]{leftmargin=1.2cm}
\onehalfspacing

\newcommand{\projectname}{E-Joutia}
\newcommand{\projectsubtitle}{Application mobile de publication et de promotion d'annonces d'occasion}
\newcommand{\screenshotplaceholder}[3]{
\begin{figure}[H]
\centering
\IfFileExists{#1}
{\includegraphics[width=0.82\textwidth]{#1}}
{
\begin{tcolorbox}[
  colback=white,
  colframe=black!30,
  width=0.82\textwidth,
  height=7.2cm,
  halign=center,
  valign=center
]
\textbf{Emplacement réservé pour la capture d'écran}\\[0.2cm]
\texttt{#1}
\end{tcolorbox}
}
\caption{#2}
\label{#3}
\end{figure}
}

\begin{document}

\begin{titlepage}
\thispagestyle{emptyclean}
\begin{center}
\vspace*{-0.8cm}

{\Large\textbf{ROYAUME DU MAROC}}\\[0.2cm]
{\normalsize Ministère de l'Enseignement Supérieur, de la Recherche Scientifique et de l'Innovation}\\[0.8cm]

\IfFileExists{images/logo_ecole.jpg}
{\includegraphics[width=0.45\textwidth]{images/logo_ecole.jpg}\\[0.8cm]}
{\vspace{0.8cm}}

{\large\textbf{Faculté des Sciences et Techniques de Tanger}}\\
{\normalsize Université Abdelmalek Essaâdi}\\[1.1cm]

\rule{\linewidth}{0.5mm}\\[0.6cm]
{\Huge\textbf{Rapport de Projet de Fin d'Études}}\\[0.6cm]
\rule{\linewidth}{0.5mm}\\[1.3cm]

{\LARGE\textbf{\projectname}}\\[0.3cm]
{\large \projectsubtitle}\\
{\normalsize avec création, validation, publication et boost d'annonces sur mobile}\\[1.2cm]

\begin{tcolorbox}[
  colback=white,
  colframe=black!55,
  boxrule=0.45pt,
  arc=0pt,
  width=0.86\textwidth,
  left=8pt,
  right=8pt,
  top=7pt,
  bottom=7pt
]
\centering
{\large\textbf{LST Ingénierie de Développement}}\\[0.18cm]
{\large\textbf{d'Applications Informatiques}}
\end{tcolorbox}

\vspace{0.65cm}

\begin{tabular}{p{0.42\textwidth} p{0.42\textwidth}}
\textbf{Réalisé par :} & \textbf{Encadré par :} \\[0.25cm]
El Bahlouli Adil & Encadrant pédagogique \\
El Boualiti Oussama & M. Serbout Zakaria \\[0.65cm]
\end{tabular}

\vfill

\begin{center}
{\small
\textbf{Période :} 01 Avril -- 01 Juin
\quad | \quad
\textbf{Année académique :} 2025--2026
\quad | \quad
\textbf{Lieu :} Tanger
}
\end{center}

\clearpage
\pagenumbering{arabic}
\setcounter{page}{1}
\end{center}
\end{titlepage}

\tableofcontents
\listoffigures
\clearpage

\chapter*{Résumé}
\addcontentsline{toc}{chapter}{Résumé}
Ce rapport présente la conception et la réalisation de \projectname{}, une application mobile dédiée à la publication d'annonces de produits d'occasion. Le projet vise à proposer une expérience fluide, moderne et adaptée aux usages mobiles, en particulier pour la création d'annonces avec photos, la validation des données saisies, la publication simulée, puis la mise en avant des annonces via une fonctionnalité de \textit{boost}.

La solution a été développée avec \textbf{React Native} et \textbf{Expo}, en utilisant \textbf{JSX} pour la structuration des interfaces. Le module \textbf{expo-image-picker} a été intégré afin de permettre l'ajout d'images depuis la galerie ou l'appareil photo. Le projet repose sur une logique front-end sans backend réel, avec simulation des traitements de publication et de sponsoring. Une attention particulière a été portée à l'ergonomie, à la cohérence visuelle et à la fidélité du rendu mobile.

Le résultat obtenu couvre l'ensemble du flux principal de publication d'annonce : création, édition, validation, publication, affichage dans \textit{Mes annonces}, sponsoring, filtrage des annonces boostées et gestion de l'état des boosts actifs ou expirés.


\chapter{Contexte général du projet}
\section{Contexte}
Le marché de l'occasion connaît une forte croissance grâce à la généralisation des usages mobiles. Les utilisateurs souhaitent aujourd'hui publier une annonce rapidement, ajouter plusieurs photos, renseigner les informations essentielles sur le produit et suivre l'état de diffusion de leur annonce depuis une interface simple et moderne.

Dans ce contexte, le projet \projectname{} a été imaginé comme une application mobile de type marketplace orientée utilisateur. L'objectif est de proposer une expérience inspirée des applications modernes de vente entre particuliers, tout en gardant une architecture simple adaptée à une démonstration académique.

\section{Problématique}
La publication d'annonces sur mobile doit répondre à plusieurs contraintes :
\begin{itemize}
  \item permettre un ajout d'images rapide depuis la caméra ou la galerie ;
  \item garantir une validation claire des champs obligatoires ;
  \item offrir un parcours utilisateur fluide jusqu'à la confirmation de publication ;
  \item différencier les annonces standard des annonces sponsorisées ;
  \item visualiser facilement les annonces boostées et les boosts expirés.
\end{itemize}

\section{Objectifs}
Les objectifs principaux du projet sont les suivants :
\begin{itemize}
  \item développer une interface mobile moderne et responsive ;
  \item implémenter un formulaire complet de publication d'annonce ;
  \item permettre l'édition d'une annonce existante ;
  \item simuler un workflow de publication sans backend réel ;
  \item intégrer une fonctionnalité de sponsoring d'annonce ;
  \item organiser les annonces selon leur état : en ligne, boostées, expirées.
\end{itemize}

\chapter{Analyse des besoins}
\section{Besoins fonctionnels}
Le système doit permettre à l'utilisateur de :
\begin{enumerate}
  \item publier une nouvelle annonce ;
  \item ajouter jusqu'à 5 photos par annonce ;
  \item supprimer une photo déjà sélectionnée ;
  \item saisir un titre, une description, un prix, une catégorie et l'état de l'objet ;
  \item valider les données avant la publication ;
  \item simuler l'envoi de l'annonce et afficher un écran de succès ;
  \item consulter la liste des annonces publiées ;
  \item modifier une annonce existante ;
  \item booster une annonce pour la sponsoriser ;
  \item afficher les boosts actifs et les boosts expirés dans des onglets dédiés.
\end{enumerate}

\section{Besoins non fonctionnels}
\begin{itemize}
  \item compatibilité avec l'environnement \textbf{Expo} ;
  \item code propre, structuré et réutilisable ;
  \item design cohérent avec les maquettes fournies ;
  \item navigation simple et compréhensible ;
  \item bonne lisibilité sur écrans mobiles ;
  \item gestion correcte des permissions caméra et galerie.
\end{itemize}

\chapter{Choix technologiques}
\section{React Native}
React Native a été retenu pour développer une application mobile multiplateforme avec une logique de composants claire et maintenable.

\section{Expo}
Expo simplifie la configuration du projet, l'exécution sur mobile et l'intégration de modules utiles pour le prototypage rapide, comme \texttt{expo-image-picker}.

\section{JSX}
Le projet a été développé exclusivement en JSX, conformément aux contraintes imposées.

\section{Expo Image Picker}
Le module \texttt{expo-image-picker} permet de sélectionner des images depuis la galerie ou de capturer des photos avec la caméra, tout en gérant les permissions nécessaires.

\chapter{Conception de l'application}
\section{Architecture générale}
L'application repose sur une architecture front-end simple organisée autour de plusieurs écrans principaux :
\begin{itemize}
  \item \texttt{CreateListingScreen.jsx} : création et modification d'annonce ;
  \item \texttt{PublishingScreen.jsx} : simulation de publication ou de boost ;
  \item \texttt{SuccessScreen.jsx} : confirmation de publication ou de boost ;
  \item \texttt{BoosterScreen.jsx} : choix de la durée de sponsoring ;
  \item \texttt{App.jsx} : orchestration des écrans, navigation locale et gestion de l'état global des annonces.
\end{itemize}

\section{Structure des données}
Chaque annonce est représentée par un objet contenant notamment :
\begin{itemize}
  \item identifiant ;
  \item titre ;
  \item description ;
  \item prix ;
  \item catégorie ;
  \item état ;
  \item tableau des photos ;
  \item date de création ;
  \item statistiques simulées (\textit{vues}, \textit{messages}, \textit{favoris}) ;
  \item état de sponsoring (\texttt{isSponsored}) ;
  \item date d'expiration du boost (\texttt{boostExpiresAt}).
\end{itemize}

\section{Flux utilisateur}
Le flux principal se déroule ainsi :
\begin{enumerate}
  \item l'utilisateur ouvre le formulaire de création ;
  \item il ajoute des photos et complète les informations produit ;
  \item le système valide les champs requis ;
  \item l'annonce est publiée via une simulation ;
  \item l'annonce apparaît dans la liste \textit{Mes annonces} ;
  \item l'utilisateur peut modifier ou booster son annonce ;
  \item en cas de boost, l'annonce est marquée comme sponsorisée puis classée dans l'onglet \textit{Booster}.
\end{enumerate}

\chapter{Réalisation}
\section{Écran de création d'annonce}
Le formulaire de publication a été conçu pour reproduire fidèlement une expérience marketplace mobile. Il comprend :
\begin{itemize}
  \item une zone de photos avec compteur ;
  \item un bouton d'ajout d'images ;
  \item les actions \textit{Prendre une photo} et \textit{Choisir depuis la galerie} ;
  \item un champ titre avec limite de 50 caractères ;
  \item un champ description détaillée ;
  \item un champ prix limité aux valeurs numériques ;
  \item la sélection de catégorie ;
  \item la sélection de l'état de l'objet ;
  \item un bouton de publication ou de modification.
\end{itemize}

\subsection{Capture d'écran du formulaire}
% Ajoutez votre image dans le dossier screenshots/, puis gardez simplement ce même nom.
\screenshotplaceholder{screenshots/create-listing-screen.png}{Écran de création d'annonce}{fig:create-listing}

\section{Validation du formulaire}
La validation appliquée couvre les règles suivantes :
\begin{itemize}
  \item titre obligatoire ;
  \item prix obligatoire et numérique ;
  \item catégorie obligatoire ;
  \item état obligatoire ;
  \item minimum une photo obligatoire.
\end{itemize}

Les erreurs sont affichées sous les champs concernés afin d'améliorer la compréhension de l'utilisateur.

\subsection{Capture d'écran de validation}
\screenshotplaceholder{screenshots/form-validation-screen.png}{Exemple de validation du formulaire}{fig:form-validation}

\section{Simulation de publication}
Le projet ne repose pas sur un backend réel. Une simulation a donc été mise en place :
\begin{itemize}
  \item construction d'un objet de type \texttt{FormData} ;
  \item affichage d'un écran de chargement ;
  \item progression visuelle avec indicateurs d'étapes ;
  \item redirection vers un écran de confirmation.
\end{itemize}

\subsection{Capture d'écran du chargement de publication}
\screenshotplaceholder{screenshots/publishing-screen.png}{Écran de chargement pendant la publication}{fig:publishing}

\section{Gestion des annonces}
La page \textit{Mes annonces} affiche les cartes d'annonces avec :
\begin{itemize}
  \item image principale ;
  \item miniatures secondaires ;
  \item titre et prix ;
  \item catégorie, état et date ;
  \item statistiques simulées ;
  \item boutons \textit{Booster} et \textit{Modifier}.
\end{itemize}

Lorsqu'une annonce possède plusieurs photos, un clic sur une miniature permet de la promouvoir comme image principale.

\subsection{Capture d'écran de la page Mes annonces}
\screenshotplaceholder{screenshots/my-listings-screen.png}{Page Mes annonces après publication}{fig:my-listings}

\section{Gestion du boost}
La fonctionnalité de sponsoring repose sur :
\begin{itemize}
  \item une page dédiée de présentation des avantages du boost ;
  \item le choix d'une durée de sponsoring ;
  \item une simulation d'activation du boost ;
  \item un écran de succès spécifique ;
  \item un badge \textit{Sponsorisée} affiché sur l'annonce ;
  \item un classement automatique dans les onglets \textit{Booster} et \textit{Expirees}.
\end{itemize}

\subsection{Capture d'écran de la page Booster}
\screenshotplaceholder{screenshots/booster-screen.png}{Écran de sélection du boost d'annonce}{fig:booster}

\subsection{Capture d'écran du succès de boost}
\screenshotplaceholder{screenshots/boost-success-screen.png}{Écran de confirmation du boost}{fig:boost-success}

\chapter{Tests et validation}
\section{Tests fonctionnels réalisés}
Les tests manuels réalisés ont permis de vérifier :
\begin{itemize}
  \item l'ouverture du formulaire de création ;
  \item l'ajout et la suppression de photos ;
  \item le respect de la limite de 5 images ;
  \item la validation des champs obligatoires ;
  \item la restriction numérique du champ prix ;
  \item la publication simulée d'une annonce ;
  \item l'affichage de l'annonce dans \textit{Mes annonces} ;
  \item la modification d'une annonce existante ;
  \item l'activation d'un boost ;
  \item l'affichage des boosts actifs et expirés.
\end{itemize}

\section{Résultats}
Les fonctionnalités principales ont été intégrées avec succès dans une expérience cohérente, proche des maquettes initiales. Le prototype obtenu est exploitable dans un environnement Expo et démontre clairement le cycle de vie d'une annonce, de sa création jusqu'à sa mise en avant sponsorisée.

\subsection{Galerie récapitulative}
\screenshotplaceholder{screenshots/success-screen.png}{Écran de succès après publication}{fig:success}

\chapter{Bilan et perspectives}
\section{Bilan}
Le projet \projectname{} a permis de mettre en pratique plusieurs compétences en développement mobile :
\begin{itemize}
  \item conception d'interfaces React Native ;
  \item gestion de l'état applicatif ;
  \item intégration de modules Expo ;
  \item validation de formulaires ;
  \item adaptation d'un design UI/UX en composants mobiles ;
  \item structuration d'un flux utilisateur complet.
\end{itemize}

\section{Limites actuelles}
Le projet reste un prototype front-end, ce qui implique certaines limites :
\begin{itemize}
  \item absence d'authentification ;
  \item absence de backend réel ;
  \item absence de persistance en base de données ;
  \item statistiques d'annonces simulées ;
  \item absence de paiement réel pour le boost.
\end{itemize}

\section{Perspectives}
Plusieurs améliorations peuvent être envisagées :
\begin{itemize}
  \item intégrer une API backend pour stocker les annonces ;
  \item ajouter l'authentification des utilisateurs ;
  \item mettre en place un système réel de paiement pour les boosts ;
  \item développer une messagerie entre acheteurs et vendeurs ;
  \item ajouter des notifications push ;
  \item connecter le projet à une base de données temps réel.
\end{itemize}

\chapter*{Conclusion}
\addcontentsline{toc}{chapter}{Conclusion}
Le projet \projectname{} a permis de concevoir une application mobile moderne de publication d'annonces d'occasion, en mettant l'accent sur l'expérience utilisateur, la cohérence visuelle et la fluidité du parcours mobile. Grâce à React Native et Expo, nous avons pu construire un prototype complet couvrant la création, l'édition, la publication et la promotion d'annonces.

Cette réalisation constitue une base solide pour une future évolution vers une plateforme marketplace complète, connectée à un backend et enrichie de services avancés.

\chapter*{Annexe A : Arborescence simplifiée du projet}
\addcontentsline{toc}{chapter}{Annexe A : Arborescence simplifiée du projet}

\begin{verbatim}
PROJET REACT NATIVE/
|-- App.jsx
|-- mock/
|   `-- options.js
|-- screens/
|   |-- BoosterScreen.jsx
|   |-- CreateListingScreen.jsx
|   |-- PublishingScreen.jsx
|   `-- SuccessScreen.jsx
|-- create-annonce-flow/
|   `-- src/routes/index.tsx
`-- package.json
\end{verbatim}


\end{itemize}
